\section{From the original isotropic averages to cuspidal theta integrals}
\label{sec:constant-terms-v33}

This chapter supplies the constant-term and cuspidality portion of
the base-theta extension in BTH12.  The finite cancellation itself is
proved in FIA1--FIA8.  We use the original full Weil datum, its actual
character, the original positive definite source, and the original
coefficient functions.  No new theta kernel is selected.  The final
holomorphic realization and the remaining smooth-growth and lowering
arguments are distinct obligations.

Write $E/F$ for the fixed CM extension, $\delta=\xi_0\in E^\times$
for its purely imaginary element, $S$ for the positive definite source,
and $T$ for the target.  Put $s=\operatorname{rank}S$,
$n=\operatorname{rank}T$, and
$W=M_{n,s}(\mathbb A_E)$.  The Hermitian Gram matrices in the original
frames are $G_S,G_T$.  We use the source's oriented form
$B=\operatorname{symplO}$, its half $B/2=\operatorname{symplOH}$,
and the actual operators
\[
 H_u f(w)=\psi\bigl(B(w,u)/2\bigr)f(w+u),\qquad
 \omega(h,g)H_u=H_{r(h,g)u}\omega(h,g).
 \tag{CTC.1}
\]
All operator identities below are on the original theta space
$\Theta$.  It consists of the smooth, finite-level, rationally
quasiperiodic functions with the prescribed character.  The Weil
datum preserves this space, preserves its Haar $L^2$ norm, has the
source's pointwise continuity, rational automorphy at zero and split
positivity.  These are the actual contracts already constructed in
CW23 and compared in BTH14; we do not assume joint orbit continuity
without proving what is needed.

% BEGIN CTC1
\paragraph{CTC1: the original transvection subgroup and its quotient.}
\label{claim:CTC1}
Let $\xi\in E^n$ be nonzero and isotropic.  Define
\[
 N_\xi=\xi\xi^*G_T,\qquad
 j(a)=\xi a\quad(a\in\mathbb A_E^s),\qquad
 R(w)=\xi^*G_Tw.
 \tag{CTC.2}
\]
Then $N_\xi^2=0$, $N_\xi^*G_T=G_TN_\xi$, and
$N_\xi\ne0$.  The last assertion follows because $G_T$ is invertible
and both factors of the rank-one matrix are nonzero.  Direct
multiplication gives
\[
 z(t)=1+(t\delta)N_\xi,\qquad
 z(t)z(t')=z(t+t'),\qquad t,t'\in\mathbb A_F.
 \tag{CTC.3}
\]
These are unitary.  Conversely, if $1+aN_\xi$ is unitary, its unitary
equation reads $(a+\bar a)G_TN_\xi=0$.  A nonzero rational entry of
$N_\xi$ is invertible in $\mathbb A_E$, so $a+\bar a=0$.
Writing $\mathbb A_E=\mathbb A_F\oplus\delta\mathbb A_F$ gives
$a=t\delta$ for a unique $t$.  Thus the generating set in the
definition of $Z_I(T,\xi)$ is already a group, and (CTC.3) is exactly
that subgroup, including the source's algebraic subgroup closure.

The map $t\mapsto z(t)$ and its inverse are continuous: recover $a$
by dividing one nonzero entry of $z-1$ by the corresponding entry of
$N_\xi$, then take its $\delta$-coordinate.  The same coordinate
description shows the image is closed.  Moreover
\[
 Z_I(T,\xi)\simeq(\mathbb A_F,+),\qquad
 Z_I(T,\xi)\cap U(T)(F)\simeq(F,+).
 \tag{CTC.4}
\]
For the rational assertion the recovered $a$ is in $E$ and its skew
condition makes $a/\delta\in F$.

Here and below the adelic vector quotients are compact.  To see this
directly, for $x\in\mathbb A_F$ choose a positive integer $q$ clearing
its finitely many finite denominators.  The Chinese remainder theorem
chooses $b\in\mathcal O_F$ with $qx-b\in q\widehat{\mathcal O}_F$.
Subtract $b/q$, then subtract an element of $\mathcal O_F$ to put the
infinite component in a fixed closed fundamental parallelepiped.
The resulting representative lies in a compact set.  The same
coordinates show that $F$ is discrete.  Products and an $F$-basis of
$E$ prove the assertion for all vector spaces used here.  Haar
measure on a quotient by a discrete cocompact subgroup has finite
positive total mass.  This can also be checked with a small open
set having disjoint rational translates and a finite cover of a
compact representative set.  Any measurable fundamental domain has
that same mass, by partitioning it and another domain by their
countably many rational translates.  In particular all the source's
transvection fundamental domains have finite, nonzero measure.

The matrix identities needed later are
\[
 R(j(a))=0,\qquad
 r(1,z(t))w=w+j(t\delta R(w)),\qquad
 r(1,z(t))j(a)=j(a),\qquad B(j(a),j(b))=0.
 \tag{CTC.5}
\]
They follow by multiplying (CTC.2), using $\xi^*G_T\xi=0$.
% END CTC1

% BEGIN CTC2
\paragraph{CTC2: a uniform pointwise bound on compact Weil orbits.}
\label{claim:CTC2}
For $f\in\Theta$ and a compact set $C$ of actual unitary pairs,
\[
                 \sup_{p\in C,\,w\in W}|\omega(p)f(w)|<\infty.
 \tag{CTC.6}
\]
Here is a proof from the source's pointwise contracts.  Take a real
basis $b_1,\ldots,b_d$ of $W_\infty$ and write $D_v$ for the
infinitesimal Heisenberg derivative.  The derivative is an ordinary
directional derivative plus a purely imaginary linear coefficient.
It preserves $\Theta$, is real-linear in $v$, and commutes with finite
Heisenberg operators.  Differentiated covariance gives
\[
 D_v\omega(p)g=\omega(p)D_{r(p)^{-1}v}g.
 \tag{CTC.7}
\]
The differentiation is justified by the $L^2$ difference-quotient
argument in WA3 and BTH5.  Iterating it uses only finitely many
derivative words of length at most $d$.

Choose a finite Heisenberg neighborhood fixing $f$.  It fixes all
these derivative words, since real and finite Heisenberg directions
commute.  Continuity of $(p,u)\mapsto r(p)^{-1}u$ and compactness of
$C$ give a single compact open finite neighborhood $K_f$ whose
translates fix every $\omega(p)D_{v_1}\cdots D_{v_a}f$, for $p\in C$
and $a\le d$.  In making this assertion, bounded $v_i$ are expanded
in the fixed basis, so only finitely many fixed words are involved.

We record the elementary estimate which converts a quotient norm
to a real-cube norm.  If $g$ has this finite level, $|g|^2$ is
constant on $K_f$-translates and periodic under $W(E)$.  For a fixed
real unit cube $Q$, integrate over $Q\times K_f$.  Its projection
to $W(E)\backslash W$ has bounded multiplicity: after choosing one
point of a fiber, all other rational differences lie in the finite
set $W(E)\cap((Q\times K_f)-(Q\times K_f))$.  The same bound holds
for every translate of this compact set.  Since $K_f$ has positive
measure, there is a constant $A$ independent of $g$ and $w$ such that
\[
 \int_{[0,1]^d}\left|g\left(w+\sum_i t_i b_i\right)\right|^2dt
             \le A\|g\|_2^2.
 \tag{CTC.8}
\]
Haar measure on $W$ factors into real Lebesgue measure and finite
Haar measure up to a fixed positive constant; that constant is
absorbed into $A$.  Translations of $|g|^2$ still have the same
rational periodicity and quotient integral, so (CTC.8) is uniform
in $w$.

For completeness the required Sobolev estimate uses only the
fundamental theorem of calculus on $[0,1]$.  Multiply a restriction
$t\mapsto g(w+tb_i)$ by a phase of absolute value one that integrates
its purely imaginary derivative coefficient.  Its ordinary
derivative then has absolute value $|D_{b_i}g(w+tb_i)|$.  Cauchy--Schwarz
and averaging the endpoint inequality give
\[
 |g(w)|^2\le 2\int_0^1|g(w+tb_i)|^2dt
                 +2\int_0^1|D_{b_i}g(w+tb_i)|^2dt.
 \tag{CTC.9}
\]
Iterate in the $d$ coordinate directions, applying the same inequality
to every derivative word which occurs.  If each cube integral of
every word of length at most $d$ is at most $M$, the result is
$|g(w)|^2\le4^dM$.  This includes $d=0$.

The vectors $r(p)^{-1}b_i$ have bounded coordinates for $p\in C$.
By (CTC.7), real multilinearity and unitarity of the original Weil
action, the $L^2$ norms of every required word applied to
$\omega(p)f$ are uniformly bounded by a finite sum of the norms
of fixed words applied to $f$.  Those norms are finite: each such
word is a continuous theta function and its absolute square descends
to the compact rational quotient.  Apply (CTC.8)--(CTC.9).  This proves
(CTC.6) without inferring joint continuity from pointwise continuity.
% END CTC2

% BEGIN CTC3
\paragraph{CTC3: local character duality and compact annihilators.}
\label{claim:CTC3}
Let $K=F_v$ at a finite place and let $\psi_v$ be the original local
character.  FIA1 proves it is nontrivial.  Every additive character
$\chi:K\to\mathbb C^\times$ trivial on an open subgroup has the form
\[
                         \chi(x)=\psi_v(ax)\quad(a\in K).
 \tag{CTC.10}
\]
We give the finite argument, rather than import a local duality
theorem.  Write $\mathfrak p^c$ for the largest fractional
$\mathcal O_K$-ideal on which $\psi_v$ is trivial.  Such $c$ exists:
continuity provides a ball on which all character values are close
to one, all integer multiples stay in that ball, and a complex
number all of whose powers stay sufficiently close to one is one.
Nontriviality excludes arbitrarily large fractional ideals from
the kernel.  Thus the kernel contains $\mathfrak p^c$ and does not
contain $\mathfrak p^{c-1}$.

Choose $m$ with $\chi(\mathfrak p^m)=1$.  For $j$ sufficiently large,
the pairing
\[
 \mathfrak p^{c-m}/\mathfrak p^{c+j}\ \times\
 \mathfrak p^{-j}/\mathfrak p^m\longrightarrow\mathbb C^\times,
                   \qquad(a,x)\longmapsto\psi_v(ax)
 \tag{CTC.11}
\]
is perfect.  Indeed, if a nonzero class of $a$ annihilates the
second factor, then $a\mathfrak p^{-j}\subset\ker\psi_v$, which
forces $v(a)-j\ge c$, a contradiction.  The two finite groups
have equal cardinalities.  Finite character orthogonality shows
their duals have those same cardinalities, proving perfection.
The set of $a\in\mathfrak p^{c-m}$ which realize $\chi$ on
$\mathfrak p^{-j}$ is therefore a nonempty compact coset of
$\mathfrak p^{c+j}$.  These cosets are nested.  Their intersection
is nonempty and realizes $\chi$ on all of $K$.  Coordinatewise the
same conclusion holds for $K^q$.

We shall also need the following consequence with its quantifiers.
For a nondegenerate $K$-bilinear pairing $b$ on two copies of $K^q$
and any neighborhood $V$ of zero in the first copy, there exists a
compact open additive subgroup $L$ of the second copy such that
\[
 \psi_v(b(\eta,L))=1\quad\Longrightarrow\quad\eta\in V.
 \tag{CTC.12}
\]
Choose linear coordinates for $b$, choose a small product of
fractional ideals contained in $V$, and take the reciprocal ideals
with conductor $\mathfrak p^c$.  Formula (CTC.11), or its valuation
argument before passing to a quotient, proves (CTC.12).  All changes
of coordinates are the original invertible Gram matrices.
% END CTC3

% BEGIN CTC4
\paragraph{CTC4: split-half coinvariant functionals separate theta vectors.}
\label{claim:CTC4}
Let $e,\bar e$ be the original pair of idempotents at a split finite
place.  Put $Y=\bar eW$ and $Y'=eW$.  If $g\in\Theta$ and
\[
 m(H_tg)=0\quad\hbox{for all }t\in Y'
 \quad\hbox{and all }m\hbox{ invariant under }H(Y),
 \tag{CTC.13}
\]
then $g=0$.  The functionals are algebraic; no continuity is imposed.
The form $B$ pairs $Y$ and $Y'$ nondegenerately, and each is isotropic.
In the original frames this follows by multiplying by $G_S^{-1}$,
$G_T$ and the nonzero $\delta$ in the trace formula for $B$.

Here is the needed algebraic separation argument.  Let a commuting
additive action $\rho$ on a complex vector space $V$ have the property
that for every $u,f$ some positive integer $N$ satisfies
$\rho(Nu)f=f$.  If $g\ne0$, there exist an additive character $\chi$
and a functional $\ell$ such that
\[
 \ell(g)\ne0,\qquad \ell(\rho(u)f)=\chi(u)\ell(f),
 \qquad \rho(u)g=g\Longrightarrow\chi(u)=1.
 \tag{CTC.14}
\]
To prove this, use the commutative group algebra
$A=\mathbb C[G]$, with basis elements $X_u$, and choose a maximal
ideal $\mathfrak m$ containing the annihilator of $g$.
Every $X_u$ has a positive power equal to one in $A/\mathfrak m$,
because it has such a power on $g$.  Thus this residue field is
algebraic over $\mathbb C$, so it is $\mathbb C$.  Write
$c:A\to\mathbb C$ for the resulting homomorphism and
$\chi(u)=c(X_u)$.

Let $V_0$ consist of those $f$ killed by some $a\in A$ with
$c(a)\ne0$.  Commutativity makes this a vector subspace, and
$g\notin V_0$.  Choose a functional on $V/V_0$ nonzero at $g$.
For given $u,f$, take a positive common multiple $N$ of a period of
$u$ on $f$ and a period on $g$.  Then $\chi(u)^N=1$.  The polynomial
\[
 q(X)=\sum_{i=0}^{N-1}X^i\chi(u)^{N-1-i}
 \quad\hbox{satisfies}\quad
 q(X)(X-\chi(u))=X^N-1,
 \quad q(\chi(u))=N\chi(u)^{N-1}\ne0.
\]
Consequently $\rho(u)f-\chi(u)f\in V_0$.  This proves every assertion
in (CTC.14), including the common eigenfunctional on all of $V$.

Apply this argument to $V=\Theta$ and $\rho(u)=H_u$, $u\in Y$.
Finite smoothness supplies the required periods: for a fixed local
vector $u$, multiplication by sufficiently large powers of the
residue characteristic puts it into any prescribed open stabilizer.
The stabilizer of $g$ is open, so (CTC.14) makes $\chi$ trivial on
an open subgroup.  CTC3 and the original nondegenerate pairing give
$t_0\in Y'$ with $\chi(u)=\psi(B(u,t_0))$.
The Heisenberg commutator is
\[
 H_uH_{t_0}=\psi(B(u,t_0))H_{t_0}H_u.
 \tag{CTC.15}
\]
It follows from (CTC.14)--(CTC.15) that
$m(f)=\ell(H_{t_0}f)$ is $H(Y)$-invariant.  But
$m(H_{-t_0}g)=\ell(g)\ne0$, contradicting (CTC.13).
An algebraic functional on $\Theta$ extends to the ambient function
space by a basis extension, exactly as required by the source.
% END CTC4

% BEGIN CTC5
\paragraph{CTC5: finite averaging transfers a small-row test to along coinvariants.}
\label{claim:CTC5}
Keep the split place of CTC4.  Suppose $g\in\Theta$ and $V$ is a
neighborhood of zero in $\mathbb A_E^s$ with the following property:
\[
 m(H_tg)=0\quad\text{whenever }t\in eW,\ R(t)\in V,
 \quad m\text{ is }H(\bar eW)\text{-invariant}.
 \tag{CTC.16}
\]
Then every functional $\ell$ invariant under all $H_{j(a)}$
satisfies $\ell(g)=0$.

The exact original trace calculation is
\[
 B\bigl(t,j(\bar s)\bigr)
   =\operatorname{Tr}_{\mathbb A_E/\mathbb A_F}
        \bigl(\delta\,(R(t)G_S^{-1})\cdot s\bigr).
 \tag{CTC.17}
\]
For $t\in eW$ this is the nondegenerate local pairing of CTC3
applied to the row coordinate $R(t)$; multiplication by
$\delta G_S^{-1}$ is invertible on that coordinate space.
Choose a compact open $L\subset e\mathbb A_E^s$ such that
triviality of (CTC.17) as a character of $L$ implies $R(t)\in V$.
Let $\alpha(s)=j(\bar s)$ for $s\in L$.  This lies in $\bar eW$
and in the along space, which is isotropic.

The continuous orbit $s\mapsto H_{\alpha(s)}g$ has an open
stabilizer in the compact group $L$, by the finite level of $g$.
That stabilizer has finite index.  For representatives $s_1,\ldots,s_q$,
where $q>0$, set $g_L=\sum_i H_{\alpha(s_i)}g$.
The sum is invariant under all $H_{\alpha(s)}$, $s\in L$.
For an $H(\bar eW)$-invariant $m$, the commutator gives
\[
 m(H_tH_{\alpha(s)}g)
          =\psi(B(t,\alpha(s)))m(H_tg).
 \tag{CTC.18}
\]
If this character of $L$ is trivial, (CTC.12), (CTC.16) and
(CTC.18) give $m(H_tg_L)=0$.  If it is nontrivial, apply (CTC.18)
to the invariant vector $g_L$ and to an $s$ where the character is
not one.  The same conclusion follows.  CTC4 now gives $g_L=0$.
Applying the along-invariant $\ell$ yields $q\ell(g)=0$, hence
$\ell(g)=0$.  This finite average is a proof about coinvariants;
it is separate from the arithmetic finite average FIA constructs
for the actual Fock vectors.
% END CTC5

% BEGIN CTC6
\paragraph{CTC6: positivity on the actual along coinvariants.}
\label{claim:CTC6}
For a transvection $z\in Z_I(T,\xi)$ supported at the split place
$e+\bar e$, there exists a single $r_z>0$ such that
\[
             \ell(\omega(1,z)f)=r_z\ell(f)
 \tag{CTC.19}
\]
for every $f\in\Theta$ and every along-invariant $\ell$.
First, the rank-one determinant formula gives
$\det(1+aN_\xi)=1+a\xi^*G_T\xi=1$.
In the original split-positivity axiom take the two determinant
lifts to be $1$.  Their splitting-character factors are exactly
$1$.  That axiom therefore supplies $r_z>0$ with (CTC.19) for
every functional invariant under the larger space $\bar eW$.
This is an application of the already constructed full Weil datum,
not an assumption that $r_z=1$.

Fix $f$ and put $g=\omega(1,z)f-r_zf$.  Write
$z^{-1}=1+aN_\xi$.  Its action on $t\in eW$ is
$r(1,z^{-1})t=t+j(aR(t))$.  If $R(t)$ is sufficiently small,
$j(aR(t))$ lies in the finite Heisenberg stabilizer of $f$.
It also lies in $eW$, as does $t$, so their Heisenberg half-phase
is zero.  Covariance consequently gives
\[
 H_t\omega(1,z)f=\omega(1,z)H_tf
                  \quad\text{when }R(t)\text{ is sufficiently small}.
 \tag{CTC.20}
\]
The split-positivity formula applied to $H_tf$ shows that $g$
satisfies (CTC.16).  CTC5 gives $\ell(g)=0$ for every along-invariant
$\ell$, proving (CTC.19).  The scalar $r_z$ came from the datum
before $f$ or $\ell$ was chosen, as required in the Fourier argument.
% END CTC6

% BEGIN CTC7
\paragraph{CTC7: a nontrivial eigenvalue for every nonzero rational row.}
\label{claim:CTC7}
Let $x\in M_{n,s}(E)$, embedded rationally in $W$, and suppose
$m=R(x)\ne0$.  Put
\[
                    N=\sum_{i,j}(G_S^{-1})_{ij}\bar m_jm_i.
 \tag{CTC.21}
\]
This is fixed by conjugation.  At any chosen complex embedding,
$G_S^{-1}$ is positive definite, so $N\ne0$.  Thus $N\in F^\times$
and $q=\delta^2N\in F^\times$.

The original host has a split place over $2$.  Use its original
split idempotent and choose a local $y\in F_v$ with
$\psi_v(y)\ne1$, by FIA1.  Let $t\in\mathbb A_F$ be supported
at this place with value $y/q$.  With $z=z(t)$ and
$b=t\delta m$, (CTC.5) gives $r(1,z)x=x+j(b)$.  The original
Hermitian trace formula, including its half, gives
\[
 h_O(j(b),x)=t\delta N,\qquad
 \tfrac12 B(j(b),x)
   =\tfrac12\operatorname{Tr}_{\mathbb A_E/\mathbb A_F}
                      (t\delta^2N)=tq.
 \tag{CTC.22}
\]
Consequently $\eta=\psi(B(j(b),x)/2)\ne1$.
No division by an integral unit at the dyadic place is asserted:
the calculation takes place in its field of fractions, and its
factor $2$ cancels before the value of the character is used.

For every along-invariant $\ell$ there is therefore a $c\ne1$
such that
\[
 \ell(H_x\omega(1,z)f)=c\ell(H_xf)
                         \qquad(f\in\Theta),\qquad c=\eta r_z.
 \tag{CTC.23}
\]
Indeed, covariance gives
$\omega(1,z)H_xf=H_{j(b)+x}\omega(1,z)f$.
Use $H_{j(b)}H_x=\eta H_{j(b)+x}$, invariance of $\ell$ under
$H_{j(b)}$, and (CTC.19) applied to $H_xf$.  This is precisely
(CTC.23).  If $c=1$, its absolute value and $r_z>0$ force
$r_z=1$, since $|\eta|=1$; then $\eta=1$, a contradiction.
The positivity of the source enters here through (CTC.21), and
the positivity of the Weil scalar enters through this last step.
% END CTC7

% BEGIN CTC8
\paragraph{CTC8: the compact-quotient Fourier averaging lemma.}
\label{claim:CTC8}
Let $Y/\Gamma$ and $Z/\Lambda$ be compact abelian quotients of
locally compact second countable groups, with finite positive Haar
measures.  Let $k(z,y)$ be bounded, separately continuous and
periodic in both variables.  Suppose a family of unitary characters
$\chi_a$ on $Y/\Gamma$ contains $1$, is closed under multiplication
and conjugation, and separates points.  Suppose also that, for each
nontrivial $\chi_a$, there are $z_a$ and $c_a\ne1$ such that
\[
 \int_{Y/\Gamma}\chi_a(y)k(z_a z,y)\,dy
       =c_a\int_{Y/\Gamma}\chi_a(y)k(z,y)\,dy
                      \quad\text{for every }z.
 \tag{CTC.24}
\]
Then $\Phi(y)=\int_{Z/\Lambda}k(z,y)\,dz$ is constant.

All integrals here are genuine integrals of bounded measurable
functions over finite measures.  Separate continuity on these
second countable spaces gives joint measurability: use countably
many shrinking neighborhoods and measurable choices of a nearby
point in one variable to express the function as a pointwise limit
of functions measurable in the product.  Dominated convergence
makes $\Phi$ continuous.  Fubini and invariance of quotient Haar
measure show that its $\chi_a$-coefficient $A_a$ satisfies
$A_a=c_aA_a$, hence $A_a=0$ for every nontrivial character in the
family.  Translation does not change the trivial coefficient.

For fixed $y_0$, the continuous function
$D(y)=\Phi(y+y_0)-\Phi(y)$ consequently has all these coefficients
zero.  To make the density step explicit, the complex linear span
of the characters is a conjugation-stable algebra containing the
constants and separating points.  Its real part is uniformly dense
in the real continuous functions by the usual elementary
Stone--Weierstrass argument: polynomial approximation to absolute
value makes its closure stable under maximum and minimum; functions
matching a prescribed continuous function at two specified points
are obtained from separation and constants; two finite-cover
arguments with maximum and minimum give a uniform approximation
to that function.  Thus its complex span uniformly approximates
$\bar D$.  Integrating $D$ against such approximations gives
$\int|D|^2=0$.  Since Haar measure is positive on nonempty open
sets and $D$ is continuous, $D=0$.  This proves the assertion for
every $y_0$, and proves constancy at every point, not merely almost
everywhere.  Normalizing either Haar measure is unnecessary.
% END CTC8

% BEGIN CTC9
\paragraph{CTC9: the actual transvection constant term is Heisenberg invariant.}
\label{claim:CTC9}
For every original Haar measure $\mu$ and rational fundamental
domain $\mathcal F$ of $Z_I(T,\xi)$, set
\[
 \mathcal C_\xi(f)=\int_{\mathcal F}\omega(1,z)f(0)\,d\mu(z).
\]
Then, for every $f\in\Theta$ and $a\in\mathbb A_E^s$,
\[
                   \mathcal C_\xi(H_{j(a)}f)=\mathcal C_\xi(f).
 \tag{CTC.25}
\]
This is the exact source theorem
\texttt{constantTerm\_heisO\_along} on its original operators.

Let $Y=\mathbb A_E^s$ and $\Gamma=j^{-1}(W(E))$.  A nonzero
coordinate of $\xi$ shows $\Gamma=E^s$: divide a rational entry of
$j(a)$ in that row by its nonzero rational coefficient.  Thus
$Y/\Gamma$ is compact by CTC1.  Isotropy and quasiperiodicity make
$a\mapsto f(j(a))$ periodic modulo $\Gamma$.  Define the functional
\[
                         \ell(f)=\int_{Y/\Gamma}f(j(a))\,da.
 \tag{CTC.26}
\]
It is well-defined and linear on $\Theta$, because every theta
function is bounded in absolute value on the compact rational
quotient of $W$.  The phase in $H_{j(b)}f(j(a))$ is one by (CTC.5);
translation invariance of quotient Haar measure therefore makes
$\ell$ invariant under every $H_{j(b)}$.

For rational matrices $x\in W(E)$ use the characters
\[
             \chi_x(a)=\psi(B(j(a),x)).
 \tag{CTC.27}
\]
They are continuous, unitary, trivial on $\Gamma$, and satisfy
$\chi_{x+x'}=\chi_x\chi_{x'}$, $\chi_{-x}=\bar\chi_x$.
They separate $Y/\Gamma$: if all their values at $a$ are one,
the rational self-annihilator calculation of CW15 gives
$j(a)\in W(E)$, hence $a\in\Gamma$.
This uses the original oriented Gram form; its rational invertible
Gram matrix identifies its dual with the rational Hom space, as
in CW15, so no change of theta model is involved.

Take
\[
                         k(z,a)=\omega(1,z)f(j(a)).
 \tag{CTC.28}
\]
It is separately continuous by the original Weil and theta
contracts.  It is periodic in $a$ as above.  It is periodic in $z$
under rational transvections: these fix $j(a)$; move $H_{j(a)}$
past them by covariance, evaluate at zero, and use rational
automorphy of the original datum.  Choose a compact representative
set for $F\backslash\mathbb A_F$ in CTC1's identification
(CTC.4).  CTC2 bounds (CTC.28) there uniformly in $a$; rational
periodicity extends the bound to every $z$.  In particular all
integrals in (CTC.25)--(CTC.28) are absolutely convergent.

If $R(x)=0$, the original trace formula makes $\chi_x=1$.
Otherwise CTC7 applies to (CTC.26).  For rational $x$ the two
half-phases in the definition and the quasiperiodicity combine to
give the full modulation formula
\[
                             H_xg(w)=\psi(B(w,x))g(w).
 \tag{CTC.29}
\]
Therefore (CTC.23), applied to $\omega(1,z)f$, is exactly
(CTC.24) for the coefficient of (CTC.28).  CTC8 shows that
$\int_{\mathcal F}k(z,a)d\mu(z)$ is independent of $a$.
Finally covariance, (CTC.5), and $H_u g(0)=g(u)$ identify that
integral with the left side of (CTC.25); at $a=0$ it is the right
side.  This proves the claimed invariance with the source's
unscaled Haar measures and with no remaining analytic theorem
hidden in the Fourier step.
% END CTC9

% BEGIN CTC10
\paragraph{CTC10: finite cancellation gives zero constant terms of the original lift.}
\label{claim:CTC10}
Let $H=U(S)(\mathbb A)$ be the original definite source group.
Assume it has the compact rational representative set provided
by the original finite class presentation in BTH3.  Let $X$ be the
original finite-dimensional coefficient space, $L:X\to\Theta$
the original linear kernel map, and $F:H\to X$ the original
continuous rationally left-invariant coefficient function.  Let
$\nu,\mathcal H$ be its original Haar measure and fundamental
domain.  Suppose that for every $x,h,g$ and rational nonzero
isotropic $\xi$, there are $q>0$ and $a_1,\ldots,a_q$ with
\[
             \sum_{i=1}^q H_{j(a_i)}\omega(h,g)L(x)=0.
 \tag{CTC.30}
\]
Then the actual convergent theta integral
\[
                 \Theta_F(g)=\int_{\mathcal H}
                                 \omega(h,g)L(F(h))(0)\,d\nu(h)
 \tag{CTC.31}
\]
has $ZCusp$ in the literal source sense: every transvection
constant term is zero, for every Haar measure, rational fundamental
domain and right translate.

First fix $x,h,g,\xi$.  Apply the linear functional
$\mathcal C_\xi$ on the finite sum (CTC.30).  CTC9 makes all
$q$ terms equal.  Since $q>0$ in $\mathbb C$,
\[
          \int_{\mathcal F}\omega(h,zg)L(x)(0)\,d\mu(z)=0.
 \tag{CTC.32}
\]
Only ordinary linearity of an integral of a finite sum is being
used; all its terms are integrable by CTC9.

For the exchange with (CTC.31), put
$K(z,h)=\omega(h,zg)L(F(h))(0)$.  Choose a basis $x_b$ of $X$.
The expansion $L(F(h))=\sum_b F_b(h)L(x_b)$ is finite.
Each $F_b$ is continuous and each fixed-vector orbit evaluation
is continuous by the Weil contract.  Thus $K$ is jointly
continuous, even though no joint-continuity axiom for the entire
infinite-dimensional theta space was assumed.  It is rationally
periodic in each group variable by the original coefficient
invariance and rational Weil automorphy.  Compact rational
representatives for $H$ and CTC1 for $Z_I$ make it bounded
globally.  Both quotient measures are finite, so its product
integral is absolutely convergent.  Fubini and (CTC.32), with
$x=F(h)$, now give
\[
 \int_{\mathcal F}\Theta_F(zg)\,d\mu(z)
 =\int_{\mathcal H}\int_{\mathcal F}
           \omega(h,zg)L(F(h))(0)\,d\mu(z)d\nu(h)=0.
 \tag{CTC.33}
\]
This proves the original definite-source lift statement in exactly
the class-presentation scope needed here.  It does not assert that
a separate unrestricted class-number theorem for every possible
definite group has been reconstructed.
% END CTC10

% BEGIN CTC11
\paragraph{CTC11: the full unipotent quotient is compact and contains the central transvections.}
\label{claim:CTC11}
Let $C\in M_{n,r}(E)$ be an original isotropic frame: $r>0$,
its columns are independent, and $C^*G_TC=0$.  The subgroup
$N_I(T,C)$ in the source is the unipotent group fixing the
isotropic space $U=\operatorname{im}C$ and acting trivially on
$U^\perp/U$ and on $T/U^\perp$.  We check this directly from
its defining equations, and prove its rational quotient compact.

Choose a rational matrix $C'$ with $C^*G_TC'=1$; this exists
because $G_T$ is nondegenerate and $C$ has independent columns.
Replacing $C'$ by $C'-\tfrac12C(C'^*G_TC')$ makes it isotropic
without changing that pairing.  Complete the resulting hyperbolic
pair to a rational orthogonal decomposition
$T=U\oplus T_0\oplus U'$.  In this basis the Gram matrix is
\[
       J=\begin{pmatrix}0&0&1\\0&J_0&0\\1&0&0\end{pmatrix},
                         \qquad J_0\text{ nondegenerate}.
 \tag{CTC.34}
\]
The source's first two equations are $(n-1)C=0$ and
$C^*G_T(n-1)=0$.  The third says that $(n-1)x\in U(\mathbb A_E)$
for every rational $x\in U^\perp$.  A rational basis of $U^\perp$
spans its adelic extension, so the third equation holds on that
whole extension.  Together with unitarity these equations say
exactly that, for an adelic matrix $A$ and skew-Hermitian $B$,
\[
 n(A,B)=
 \begin{pmatrix}
  1&-A^*J_0&B-\tfrac12 A^*J_0A\\
  0&1&A\\0&0&1
 \end{pmatrix},\qquad B^*=-B.
 \tag{CTC.35}
\]
For example the first two equations fix the first block column
and last block row; the third fixes the middle diagonal block;
the remaining unitary equations determine the upper-middle
block and the Hermitian part of the upper-right block.  This
derives (CTC.35), including when $T_0$ has dimension zero.

Direct multiplication gives
\[
 n(A,B)n(A',B')=
 n\bigl(A+A',B+B'+\tfrac12(A'^*J_0A-A^*J_0A')\bigr).
 \tag{CTC.36}
\]
Thus these matrices already form a closed subgroup.  The source's
subgroup closure adds nothing.  The kernel of $n(A,B)\mapsto A$
is the central additive space of skew-Hermitian $B$; the quotient
is the additive $E$-matrix space of $A$.  Both are adelic vector
spaces defined over $F$, and their rational points are their
original rational subspaces.

For any $n(A,B)$, left multiply by a rational $n(a,0)$ to put
$A+a$ into a fixed compact vector representative set.  Then left
multiply by a rational central $n(0,b)$ to put its new central
coordinate into a fixed compact vector representative set.
CTC1 supplies both compact sets.  Their image under (CTC.35)
is compact and meets every left rational coset.  This proves
\[
                     N_I(T,C)(F)\backslash N_I(T,C)(\mathbb A)
                              \quad\text{is compact}.
 \tag{CTC.37}
\]
The rational subgroup is discrete in these coordinates.
Product additive Haar measure in $(A,B)$ is left and right
invariant by the shear law (CTC.36).  Every rational fundamental
domain therefore has finite positive measure, and every continuous
rationally invariant function has an absolutely convergent
integral over it.

Finally choose any column $\xi=C_j$.  It is nonzero and isotropic,
and
$N_\xi=C E_{jj}C^*G_T$.
For every defining matrix $n$ of $N_I$, the equations above give
$(n-1)N_\xi=N_\xi(n-1)=0$.
Consequently $Z_I(T,\xi)\subset N_I(T,C)$ and this subgroup is
central.  These are the exact source inclusions, not an appeal
to an unspecified parabolic conjugacy criterion.
% END CTC11

% BEGIN CTC12
\paragraph{CTC12: zero transvection constants imply genuine cuspidality.}
\label{claim:CTC12}
Let $\varphi$ be continuous and rationally left-invariant on the
original target group, and assume $ZCusp(T,\varphi)$.  Then
$IsCuspFun(T,\varphi)$ holds with genuinely convergent constant
terms.  The finite-level part of $IsAutFun$ is not needed for
this implication.

Fix an original isotropic frame $C$ and a right translate $g$.
Let $N=N_I(T,C)$, $\Gamma=N\cap U(T)(F)$ and
$F_g(n)=\varphi(ng)$.  CTC11 makes $F_g$ continuous on a compact
rational quotient, hence bounded and integrable for every
original Haar fundamental domain $\mathcal H_N$.  Choose a
column $\xi$ of $C$ and a rational fundamental domain
$\mathcal F_Z$ of the central subgroup $Z=Z_I(T,\xi)$.
Its measure $m_Z$ is positive and finite by CTC1.

For each $z\in Z$, centrality makes $n\mapsto zn$ descend to
the left quotient $\Gamma\backslash N$.  It preserves its
Haar measure.  One can verify this directly: $z^{-1}\mathcal H_N$
is another rational fundamental domain, and partitioning the two
domains into rational translates gives equality of their integrals.
Thus
\[
 \int_{\mathcal H_N}F_g(zn)\,dn
                   =\int_{\mathcal H_N}F_g(n)\,dn.
 \tag{CTC.38}
\]
The same equality for $|F_g|$ and finiteness of $m_Z$ justify
Fubini over $\mathcal F_Z\times\mathcal H_N$.  But $ZCusp$
says $\int_{\mathcal F_Z}F_g(zn)\,dz=0$ for every $n$.
Integrating and using (CTC.38) gives
\[
                 0=m_Z\int_{\mathcal H_N}\varphi(ng)\,dn.
 \tag{CTC.39}
\]
Since $m_Z>0$, the integral is zero.  The source quantifies
over all such frames, measures, domains and $g$; all have been
retained.  If no isotropic frame exists, the condition is vacuous.

The source's abstract measure lemma also covers nonintegrable
functions because its total Bochner integral is defined to be
zero there.  That convention is not needed in this human proof:
CTC11 proves integrability of every continuous automorphic
function on every unipotent quotient appearing in the definition.
% END CTC12

% BEGIN CTC13
\paragraph{CTC13: the exact base-kernel cuspidality output and its remaining interfaces.}
\label{claim:CTC13}
Fix one original base index, parameter and source coefficient
vector in BTH12.  Use its original polynomial kernel map $L$,
the same full Weil datum, and the original finite class presentation.
The finite-average endpoint of FIA applies to every translated
polynomial kernel in (CTC.30), with a positive period and one
averaging list independent of the polynomial.  Linearity therefore
also gives (CTC.30) on the entire original finite coefficient
space; equivalently one may use FIA's universal polynomial
statement directly for each $x$.  BTH3 gives convergence of
(CTC.31) and its compact source representative set.
CTC10 now gives $ZCusp$ of the actual scalar lift.

For clarity, the two hypotheses needed to apply CTC12 can also
be checked here without waiting for the full automorphic
membership theorem.  Rational target invariance follows from
the actual rational Weil identity evaluated at zero, with the
source variable unchanged.  For continuity near $g_0$, restrict
$g$ to a compact neighborhood, choose a finite basis of the
coefficient space, and restrict $h$ to the compact source
representative set.  The finite-basis expansion used in CTC10
is jointly continuous on this compact set, hence uniformly
bounded; it is rationally invariant in $h$.  Dominated convergence
for the finite source quotient measure makes $\Theta_F(g)$
continuous at $g_0$.  Therefore CTC12 proves $IsCuspFun$ of
the original lift.

This supplies the literal isotropic-average, zero-constant-term
and cuspidality portion of the source's
\texttt{ExistsThetaLiftLinearMapCuspFormMultiplicity} construction.
It applies for each original base parameter and uses no independently
selected replacement kernel or scalar normalization.  It does
\emph{not} by itself construct the original smooth cuspidal
multiplicity-space extension. That still needs the original
smoothness and uniform moderate-growth
derivative estimates, lowering and prescribed lowest-type conditions,
the holomorphic realization into the actual Hilbert multiplicity
model. Separately, the original source route still needs
$\mathrm{PIoP}$: one common $\mathrm{ioP}$ satisfying
$\mathrm{IsIoW}$ for every original parameter. This is not the
smooth extension itself. The same-Bundle readings and rational
test identities recorded in BTH12 also remain. Those obligations
must be supplied with their original universal parameter
quantifiers. The present result
removes the cuspidality subargument from that residual task; it
does not remove the entire producer leaf.
% END CTC13


